\chapter{线性方程组}

如何求解$n$元线性方程组？
尽管Cramer法则给出了解的公式，但其算法的时间复杂度$\bigo(n \cdot (n + 1)!)$是不可接受的。

\iffalse
求解线性代数方程组就相当于求解矩阵式$Ax=b$：
\begin{equation}
    \begin{bmatrix}
        a_{11}&\cdots&a_{1n}\\
        \vdots&\ddots&\vdots\\
        a_{n1}&\cdots&a_{nn}
    \end{bmatrix}\begin{bmatrix}
        x_1\\\vdots\\x_n
    \end{bmatrix}=\begin{bmatrix}
        b_1\\\vdots\\b_n
    \end{bmatrix}.
\end{equation}
因此有必要了解一些对矩阵的操作。

\section{矩阵操作}

\begin{definition}
    {稀疏矩阵}{sparse matrix}
    如果一个矩阵绝大多数元素是0，则称其是稀疏的(sparse)。
\end{definition}

\begin{definition}
    {秩一矩阵}{rank-1 matrix}
    若一个矩阵可以表示成$A=uv\adjoint$，则其秩为一。
\end{definition}

\begin{theorem}
    {奇异值分解}{singular value}
    根据奇异值分解，可以将矩阵表示成秩一矩阵的线性组合：
    \begin{equation}
        A=U\Sigma V\adjoint=\sum_i\sigma_iu_iv_i\adjoint,
    \end{equation}
    其中$U,V$为幺正矩阵。
\end{theorem}

\begin{theorem}
    {矩阵的Hierarchical表示}{Hierarchical matrix}
    考虑分块矩阵
    \[
        A=\begin{bmatrix}
            A_{11}&A_{12}\\
            A_{21}&A_{22}
        \end{bmatrix}
    \]
    $A$一般不是稀疏的，但非对角元$A_{12},A_{22}$是稀疏的，则$Ax$可以分块地写成：
    \[
        Ax=\begin{bmatrix}
            A_{11}&A_{12}\\
            A_{21}&A_{22}
        \end{bmatrix}\begin{bmatrix}
            x_1\\x_2
        \end{bmatrix}=\begin{bmatrix}
            A_{11}x_1+A_{12}x_2\\
            A_{21}x_1+A_{22}x_2
        \end{bmatrix},
    \]
    对$A_{11}x,A_{22}x_2$递归处理。
    \begin{notecaptionblock}
        \tikzfigure{hierarchical}
        \captionof{figure}{Hierarchical算法示意图，白块表示稀疏部分}
        \label{fig:Hierarchical}
    \end{notecaptionblock}
\end{theorem}

\begin{theorem}
    {矩阵乘法的Strassen算法}{}
    对矩阵乘法$C=AB$分块得：
    \[
        \begin{bmatrix}
            C_{11}&C_{12}\\
            C_{21}&C_{22}
        \end{bmatrix}=\begin{bmatrix}
            A_{11}&A_{12}\\
            A_{21}&A_{22}
        \end{bmatrix}\begin{bmatrix}
            B_{11}&B_{12}\\
            B_{21}&B_{22}
        \end{bmatrix},
    \]
    定义 
    \begin{subequations}
        \begin{align}
            M_1&=(A_{11}+A_{22})(B_{11}+B_{22}),\\
            M_2&=(A_{21}+A_{22})B_{11},\\
            M_3&=A_{11}(B_{12}-B_{22}),\\
            M_4&=A_{22}(B_{21}-B_{11}),\\
            M_5&=(A_{11}+A_{12})B_{22},\\
            M_6&=(A_{21}-A_{11})(B_{11}+B_{12}),\\
            M_7&=(A_{12}-A_{22})(B_{21}+B_{22}),
        \end{align}
    \end{subequations}
    则
    \begin{subequations}
        \begin{align}
            C_{11}&=M_1+M_4-M_5+M_7,\\
            C_{12}&=M_3+M_5,\\
            C_{21}&=M_2+M_4,\\
            C_{22}&=M_1-M_2+M_3+M_6.
        \end{align}
    \end{subequations}
    Strassen算法将分块矩阵的乘法从直接法的8次降低到了7次，由此分而治之，矩阵乘法的时间复杂度便从$\bigo(n^3)$降低到了$\bigo(n^{\log_27})=\bigo(n^{2.807})$。
\end{theorem}

\begin{definition}
    {离散Fourier变换}{discrete Fourier transform}
    式\eqref{eqn:Fourier betaj}中定义了一个线性变换，称为离散Fourier变换(discrete Fourier transform, DFT)
    \[
        X_n=\sum_{m=0}^{N-1}x_m\omega^{mn},
    \]
    其中$\omega:=\e^{-\i2\pi/N}$是$N$次单位根，对应的变换矩阵为
    \begin{eqnarray}
        F=\begin{bmatrix}
            1&1&1&\cdots&1\\
            1&\omega&\omega^2&\cdots&\omega^{N-1}\\
            1&\omega^2&(\omega^2)^2&\cdots&(\omega^2)^{N-1}\\
            \vdots&\vdots&\vdots&\ddots&\vdots&\\
            1&\omega^{N-1}&(\omega^{N-1})^2&\cdots&(\omega^{N-1})^{N-1}
        \end{bmatrix}
    \end{eqnarray}
    事实上$F$上只有$n$个不同的元素。其逆变换为
    \begin{equation}
        F\inv=\frac1N\bar F.
    \end{equation}
\end{definition}

\begin{theorem}
    {Cooley--Tukey快速Fourier变换}{Cooley-Tukey fast Fourier transform}
    以基2 (radix-2)的情形为例，即$N=2^M$。
    将$X_n$的求和分成偶数项$E_n$和奇数项$O_n$
    \begin{align*}
        X_n&=\sum_{k=0}^{N/2-1}x_{2k}\omega^{2kn}+\sum_{k=0}^{N/2-1}x_{2k+1}\omega^{(2k+1)n}\\
        &=\sum_{k=0}^{N/2-1}x_{2k}\omega^{2kn}+\omega^n\sum_{k=0}^{N/2-1}x_{2k+1}\omega^{2kn}=:E_n+\omega^nO_n,
    \end{align*}
    由于$\omega^N=1$，注意到
    \begin{align*}
        X_{n+N/2}&=\sum_{k=0}^{N/2-1}x_{2k}\omega^{2k(n+N/2)}+\sum_{k=0}^{N/2-1}x_{2k+1}\omega^{(2k+1)(n+N/2)}\\
        &=\sum_{k=0}^{N/2-1}x_{2k}\omega^{2kn}-\omega^n\sum_{k=0}^{N/2-1}x_{2k+1}\omega^{2k}=E_n-\omega^nO_n,
    \end{align*}
    由此便将$N$个$X_n$求和($N^2$)转化成了$N/2$个$E_n,O_n$求和($N^2/2$)。
    采用分而治之的算法思想，可以将DFT的时间复杂度从矩阵向量乘法的$\bigo(N^2)$优化到$\bigo(N\log N)$，这称为快速Fourier变换(fast Fourier transform, FFT)。
    \begin{notecaptionblock}
        \tikzfigure[width=0.8\linewidth]{fft}
        \captionof{figure}{FFT算法示意图($N=2^3=8$)，\href{https://tikz.net/fast-fourier-transform/}{作者：Alexandros Tsagkaropoulos}}
        \label{fig:FFT}
    \end{notecaptionblock}
    \tcblower
    对于非基2的情形，$\omega^k$的周期不是基2的，做处理：
    \[
        X_n=\sum_{m=0}^{N-1}x_m\omega^{-mn}=\sum_{m=0}^Nx_m\omega^{[(m-n)^2-m^2-n^2]/2}.
    \]
    定义$\nu_k:=\omega^{k^2/2}$，记$Y_n:=\nu_nX_n,\;z_m:=\nu_m\inv x_m$，则有卷积形式：
    \[
        Y_n=\sum_{m=0}^{N-1}z_m\nu_{m-n}.
    \]
    可以用0将$z_m,\nu_m$延拓，使其周期是一个比$N$大的基2数$N'$。
    再在两端做DFT：
    \begin{align*}
        \sum_{n=0}^{N'-1}Y_n\omega_{N'}^{nk}&=\sum_{n=0}^{N'-1}\sum_{m=0}^{N'-1}z_m\nu_{m-n}\omega_{N'}^{nk}\\
        &=\sum_{m=0}^{N'-1}z_m\omega_{N'}^{mk}\sum_{n=0}^{N'-1}\nu_{m-n}\omega_{N'}^{(n-m)k}=\sum_{m=0}^{N'-1}z_m\omega_{N'}^{mk}\sum_{m=0}^{N'-1}\nu_m\omega_{N'}^{mk}
    \end{align*}
    因此通过对$z_m,\nu_m$做两次DFT、一次向量分量积、一次逆DFT便可得到$Y_n$。
\end{theorem}

\begin{theorem}
    {周期Toeplitz变换}{periodic Toeplitz transform}
    $n$阶矩阵$A$若满足$a_{ij}=c_{i-j}$且序列$c_k$周期为$n$，则$A$称为周期Toeplitz矩阵，且
    \begin{equation}
        A=F\inv\Lambda F,
    \end{equation}
    其中$\Lambda=\diag(\lambda_0,\ldots,\lambda_{n-1})$且
    \begin{equation}
        \lambda_i=\sum_{j=0}^{n-1}c_j\omega^{-ij}.
    \end{equation}
\end{theorem}

\begin{remark}
    $n$阶矩阵$A$若满足$a_{ij}=c_{i-j}$，则该矩阵可以扩展成$2n$阶的周期Toeplitz矩阵。
    \[
        \begin{bmatrix}
            A&*\\ *&*
        \end{bmatrix}
    \]
\end{remark}
\fi

\section{Gauss消元法}
\label{sec:Gauss elimination}

% \paragraph{基本思路}

% 将$Ax=b$的求解问题转化为一个与之等价的容易求解的问题$\tilde A\tilde x=\tilde b$。

\subsection{Gauss消元法}



\begin{theorem}
    {Gauss消元法}{Gauss elimination}
    线性方程组形如
    \begin{equation*}
        \eqnset{\begin{aligned}
            a_{11}^{(1)}x_1 + a_{12}^{(1)}x_2 + \cdots + a_{1n}^{(1)}x_n &= b_1^{(1)}\\
            a_{21}^{(1)}x_1 + a_{22}^{(1)}x_2 + \cdots + a_{2n}^{(1)}x_n &= b_2^{(1)}\\
            &\vdots\\
            a_{n1}^{(1)}x_1 + a_{n2}^{(1)}x_2 + \cdots + a_{nn}^{(1)}x_n &= b_n^{(1)}
        \end{aligned}}
    \end{equation*}
    如果$a_{11}^{(1)} \neq 0$，可将第一行的$-a_{i1}^{(1)}/a_{11}^{(1)}$倍加到第$i$行($i = 2,3,\ldots,n$)，得到一个等价方程组
    \begin{equation*}
        \eqnset{\begin{aligned}
            a_{11}^{(1)}x_1 + a_{12}^{(1)}x_2 + \cdots + a_{1n}^{(1)}x_n &= b_1^{(1)}\\
            a_{22}^{(2)}x_2 + \cdots + a_{2n}^{(2)}x_n &= b_2^{(2)}\\
            &\vdots\\
            a_{n2}^{(2)}x_2 + \cdots + a_{nn}^{(2)}x_n &= b_n^{(2)}
        \end{aligned}}
    \end{equation*}
    如果$a_{22}^{(2)} \neq 0$，便可以此类推……最终得到一个等价的上三角线性方程组：
    \begin{equation*}
        \eqnset{\begin{aligned}
            a_{11}^{(1)}x_1 + a_{12}^{(1)}x_2 + \cdots + a_{1n}^{(1)}x_n &= b_1^{(1)}\\
            a_{22}^{(2)}x_2 + \cdots + a_{2n}^{(2)}x_n &= b_2^{(2)}\\
            &\vdots\\
            a_{nn}^{(n)}x_n &= b_n^{(n)}
        \end{aligned}}
    \end{equation*}
    便不难从下至上地解得：$x_n = b_n^{(n)}/a_{nn}^{(n)}$，
    \begin{equation}
        x_i = \frac{1}{a_{ii}^{(i)}}\biggp{b_{i}^{(i)} - \sum_{j=i+1}^na_{ij}^{(i)}x_j},\quad i = n - 1,\ldots,2,1.
    \end{equation}
\end{theorem}

\begin{corollary}
    Gauss消元法的复杂度为$\bigo(n^3)$。
\end{corollary}

\begin{remark}
    从上面的算法过程中可见，
    一旦第$k$步$a_{kk}^{(k)} = 0$，(顺序) Gauss消元法就不能继续进行下去。
    但是当第$k + 1,\ldots,n$行中存在$a_{ik}^{k} \neq 0$，就可以交换$k,i$行，从而使算法继续。
\end{remark}

\begin{remark}
    此外，即使$a_{kk}^{(k)} \neq 0$但$\abs{a_{kk}^{(k)}} \ll 1$，也会出现大数除小数导致精度下降的问题。
\end{remark}

\begin{theorem}
    {列主元方法}{pivoting technique}
    在第$k$步消去之前，找到绝对值最大的主元(pivot)：
    \begin{equation}
        i_k = \argmax_{k\leq i\leq n}\abs{a_{ik}^{(k)}},
    \end{equation}
    然后交换第$k,i_k$行。
\end{theorem}

\subsection{LU分解}

下面我们从矩阵角度考察Gauss消元法。
线性方程组等价于矩阵方程
\[
    \begin{bmatrix}
        a_{11}&\cdots&a_{1n}\\
        \vdots&\ddots&\vdots\\
        a_{n1}&\cdots&a_{nn}
    \end{bmatrix}\begin{bmatrix}
        x_1\\\vdots\\x_n
    \end{bmatrix} = \begin{bmatrix}
        b_1\\\vdots\\b_n
    \end{bmatrix}.
\]

\begin{definition}
    {初等矩阵}{elementary matrix}
    给定(实)向量$u,v$和标量$\sigma$，形如
    \begin{equation}
        E = I - \sigma uv\transp
    \end{equation}
    的称为(实)初等矩阵(elementary matrix)。
\end{definition}

\begin{corollary}
    初等矩阵的逆也是同类型的初等矩阵：
    \begin{equation}
        (I - \sigma uv\transp)\inv = I - \frac{\sigma}{\sigma v\transp u - 1}uv\transp.
    \end{equation}
\end{corollary}

\begin{example}
    {}{}
    \begin{itemize}
        \item 初等排列矩阵：
        \begin{equation}\label{eq:na-elementary-permutation}
            P_{ij} = I - (e_i - e_j)(e_i - e_j)\transp = I - e_{ii} - e_{jj} + e_{ij} + e_{ji}.
        \end{equation}
        左乘$P_{ij}$即互换第$i,j$行，右乘即互换第$i,j$列；
        \item 倍加矩阵：
        \begin{equation}\label{eq:na-elementary-row-addition}
            I + \alpha e_ie_j\transp = I + \alpha e_{ij},
        \end{equation}
        左乘$I + \alpha e_{ij}$即将第$i$行的$\alpha$倍加到第$j$行上。
    \end{itemize}
\end{example}

\begin{definition}
    {初等单位下三角矩阵}{elementary lower-triangle matrix}
    形如
    \begin{equation}
        L_i = I + \ell_ie_i\transp = \begin{bmatrix}
            1\\ &\ddots\\ &&1\\ &&\ell_{i+1,i}&1\\ &&\vdots&&\ddots\\ &&\ell_{ni}&&&1
        \end{bmatrix}
    \end{equation}
    称为第$i$列的初等单位下三角矩阵。其中向量$\ell_i$的前$i$个分量为0。
\end{definition}

\begin{corollary}
    初等单位下三角矩阵的逆也是初等单位下三角矩阵
    \begin{equation}
        (I + \ell_ie_i\transp)\inv = I - \ell_ie_i\transp.
    \end{equation}
\end{corollary}

\begin{corollary}
    对角元均为1的下三角矩阵称为单位下三角矩阵，可以写成
    \[
        L = L_1L_2\cdots L_{n-1} = \begin{bmatrix}
            1\\\ell_{21}&1\\ \vdots&\vdots&\ddots\\
            \ell_{n1}&\ell_{n2}&\cdots&1
        \end{bmatrix}.
    \]
\end{corollary}

\begin{theorem}
    {矩阵的LU分解}{LU decomposition exists}
    根据Gauss消元法的过程，记
    \begin{align*}
        A^{(k)} &:= \begin{bmatrix}
            a_{11}^{(1)}&\cdots&a_{1k}^{(1)}&\cdots&a_{1n}^{(1)}\\
            &\ddots&\vdots&\ddots&\vdots\\
            &&a_{kk}^{(k)}&\cdots&a_{kn}^{(k)}\\
            &&\vdots&\ddots&\vdots\\
            &&a_{nk}^{(k)}&\cdots&a_{nn}^{(k)}
        \end{bmatrix},\\
        b^{(k)} &:= \begin{bmatrix}
            b_1^{(1)}\\\vdots\\b_k^{(k)}\\\vdots\\b_n^{(k)}
        \end{bmatrix},\\
        \ell_k &:= -\frac1{a_{kk}^{(k)}}\begin{bmatrix}
            0\\\vdots\\0\\a_{k+1,k}^{(k)}\\\vdots\\a_{nk}^{(k)}
        \end{bmatrix}.
    \end{align*}
    相应的初等单位下三角矩阵为$L_k = I + \ell_ke_k\transp$。
    增广矩阵有递推关系：
    \begin{equation}
        [A^{(k+1)}\enspace b^{(k+1)}] = L_k[A^{(k)}\enspace b^{(k)}],
    \end{equation}
    即
    \[
        A^{(n)} = L_{n-1}\cdots L_1A^{(1)} \iff A^{(1)} = L_1\inv\cdots L_{n-1}\inv A^{(n)},
    \]
    则$L = L_1\inv\cdots L_{n-1}\inv$为单位下三角矩阵，$U = A^{(n)}$为上三角矩阵，这样就将$A^{(1)}$分解成了下上三角矩阵的乘积。
\end{theorem}

\begin{remark}
    根据Gauss消元法的过程可见，LU分解的前提是$a_{11}^{(1)},\ldots,a_{nn}^{(n)}$均不为0。
\end{remark}

\begin{theorem}
    {三角分解定理}{LU decomposition}
    给定矩阵$A$，若其顺序主子式
    \begin{equation}
        \Delta_i := \begin{vmatrix}
            a_{11}&\cdots&a_{1i}\\
            \vdots&\ddots&\vdots\\
            a_{i1}&\cdots&a_{ii}
        \end{vmatrix},\quad i = 1,\ldots,n
    \end{equation}
    均不为0，则存在唯一的单位下三角矩阵$L$和上三角矩阵$U$使得$A = LU$。
\end{theorem}
\begin{proof}
    通过数学归纳法，可证明：$a_{11}^{(1)},\ldots,a_{ii}^{(i)} \neq 0 \iff \Delta_1,\ldots,\Delta_i \neq 0$，此时
    \begin{equation}
        \Delta_i = a_{11}^{(1)}\cdots a_{ii}^{(i)},
    \end{equation}
    若存在$L_1,U_1$和$L_2,U_2$使得
    \[
        A = L_1U_1 = L_2U_2,
    \]
    两边左乘$L_1\inv$，右乘$U_2\inv$得到：
    \[
        U_1U_2\inv = L_1\inv L_2,
    \]
    此式左边为上三角矩阵，右边为单位下三角矩阵，故只能是单位矩阵$I$，即$U_1 = U_2,L_1 = L_2$。
\end{proof}

\begin{example}
    {LU分解的例子}{}
    \[
        \begin{bmatrix}
            2&1&1&0\\4&3&3&1\\8&7&9&5\\6&7&9&8
        \end{bmatrix} = \begin{bmatrix}
            1\\2&1\\4&3&1\\3&4&1&1
        \end{bmatrix}\begin{bmatrix}
            2&1&1&0\\ &1&1&1\\ &&2&2\\ &&&2
        \end{bmatrix}.
    \]
\end{example}

\begin{remark}
    $L$有$n(n - 1)/2$个变量，$U$有$n(n + 1)/2$个变量，可以直接由$A$确定$L,U$。
\end{remark}

\begin{theorem}
    {直接LU分解法(Doolittle分解法)}{Doolittle decomposition}
    写成分块矩阵的形式：
    \[
        L^{(k)} = \begin{bmatrix}
            1\\\ell_{n-k+1}&L^{(k-1)}
        \end{bmatrix},\quad
        U^{(k)} = \begin{bmatrix}
            u_{u-k+1,u-k+1}&u_{u-k+1}\transp\\ &U^{(k-1)}
        \end{bmatrix},
    \]
    则
    \[
        A^{(k)} = L^{(k)}U^{(k)} = \begin{bmatrix}
            u_{u-k+1,u-k+1}&u_{u-k+1}\transp\\
            u_{u-k+1,u-k+1}\ell_{n-k+1}&A^{(k-1)} + \ell_{n-k+1}u_{u-k+1}\transp
        \end{bmatrix}.
    \]
    由此可确定$u,\ell$，同时$A^{(k)}$的阶数减1.
\end{theorem}

\begin{theorem}
    {Thomas方法}{Thomas technique}
    对于三对角矩阵有如下LU分解：
    \begin{equation}\label{eq:na-tridiagonal-lu}
        \begin{bmatrix}
        b_1&c_1\\a_2&b_2&\ddots\\ &\ddots&\ddots&c_{n-1}\\ &&a_n&b_n
        \end{bmatrix} = \begin{bmatrix}
            1\\\ell_2&1\\&\ddots&\ddots\\ &&\ell_n&1
        \end{bmatrix}\begin{bmatrix}
            u_1&c_1\\ &u_2&\ddots\\ &&\ddots&c_{n-1}\\ &&&u_n
        \end{bmatrix},
    \end{equation}
    可直接乘开得到
    \begin{subequations}
        \begin{align}
            u_1 &= b_1,\\
            \ell_i &= \frac{a_i}{u_{i-1}},\\
            u_i &= b_i - \ell_ic_{i-1}.
        \end{align}
    \end{subequations}
\end{theorem}
    
三角分解\thmref{thm:LU decomposition} 给出了LU分解的条件。
而对于一般的可逆矩阵，也可以通过换行实现LU分解。

\begin{theorem}
    {一般三角分解定理}{LUP decomposition}
    若$A$可逆，则存在排列矩阵$P$、单位下三角矩阵$L$和上三角矩阵$U$使得
    \begin{equation}
        PA = LU.
    \end{equation}
\end{theorem}

\begin{proof}
    考虑列主元方法的Gauss消元法，第$k$步交换$k,i_k$行，则
    \[
        A^{(k+1)} = L_kI_{ki_k}A^{(k)},
    \]
    即
    \[
        A^{(n)} = L_{n-1}I_{n-1,i_{n-1}}\cdots L_1I_{1i_1}A^{(1)},
        \iff
        A^{(1)} = I_{1i_1}L_1\inv\cdots I_{n-1,i_{n-1}}L_{n-1}\inv A^{(n)},
    \]
    定义
    \[
        P_k = I_{n-1,i_{n-1}}\cdots I_{ki_k},
    \]
    则$P_k\transp P_{k+1} = I_{ki_k}$，进而
    \[
        P_1A^{(1)} = P_2L_1\inv P_2\transp P_3L_2\inv\cdots P_{n-1}L_{n-1}\inv A^{(n)},
    \]
    易得$P_{k+1}e_k = e_k$，故
    \[
        L_k' := P_{k+1}L_k\inv P_{k+1}\transp = P_{k+1}(I - \ell_ke_k\transp)P_{k+1}\transp = I - P_{k+1}\ell_ke_k\transp,
    \]
    仍然是第$k$列的初等单位上三角矩阵，
    令$P = P_1,\;L = L_1'\cdots L_{n-2}'L_{n-1},\;U = A^{(n)}$即得。
\end{proof}

\subsection{Cholesky分解}

下面再看对称矩阵的三角分解。

\begin{theorem}
    {Cholesky分解}{Cholesky decomposition}
    若$A$实对称正定，则存在唯一的对角元素为正的下三角矩阵$L$使得
    \begin{equation}
        A = LL\transp.
    \end{equation}
\end{theorem}

\begin{proof}
    采用加边Cholesky分解法：写成分块矩阵的形式
    \[
        L_i = \begin{bmatrix}
            L_{i-1}\\\ell_{i-1}\transp&\ell_{ii}
        \end{bmatrix},\quad A_i = \begin{bmatrix}
            A_{i-1}&a_{i-1}\\ a_{i-1}\transp&a_{ii}
        \end{bmatrix}
    \]
    满足$A_i = L_iL_i\transp$，
    可得 
    \begin{subequations}
        \begin{align}
            \ell_{i-1} &= L_{i-1}\inv a_{i-1},\\
            \ell_{ii} &= \sqrt{a_{ii} - \ell_{i-1}\transp\ell_{i-1}}.
        \end{align}
    \end{subequations}
    从$\ell_{11} = \sqrt{a_{11}}$出发便可迭代得到整个$L$。
\end{proof}

\begin{remark}
    这种算法特别适合稀疏矩阵。
\end{remark}

\subsection{稳定性分析}
\label{sec:stability analysis}

在用直接法求解$Ax = b$的过程中，由于舍入误差的存在，必然会导致结果产生误差。因而有必要对可能产生的误差作一估计。
% 通常我们假设在数值处理的过程中计算都是精确的.
\begin{example}
    {数据的微小变化导致解的巨大变化}{}
    方程组
    \[
        \begin{bmatrix}
            10&7&8&7\\
            7&5&6&5\\
            8&6&10&9\\
            7&5&9&10
        \end{bmatrix}\begin{bmatrix}
            x_1\\x_2\\x_3\\x_4
        \end{bmatrix} = \begin{bmatrix}
            32\\23\\33\\31
        \end{bmatrix} \implies \begin{bmatrix}
            x_1\\x_2\\x_3\\x_4
        \end{bmatrix} = \begin{bmatrix}
            1\\1\\1\\1
        \end{bmatrix}.
    \]
    对向量$b$数据做微小的修改
    \[
        \begin{bmatrix}
            10&7&8&7\\
            7&5&6&5\\
            8&6&10&9\\
            7&5&9&10
        \end{bmatrix}\begin{bmatrix}
            x_1'\\x_2'\\x_3'\\x_4'
        \end{bmatrix} = \begin{bmatrix}
            32.1\\22.9\\33.1\\30.9
        \end{bmatrix} \implies \begin{bmatrix}
            x_1'\\x_2'\\x_3'\\x_4'
        \end{bmatrix} = \begin{bmatrix}
            9.2\\-12.6\\4.5\\-1.1
        \end{bmatrix}.
    \]
    对系数矩阵$A$数据做微小的修改
    \[
        \begin{bmatrix}
            10&7&8.1&7.2\\
            7.08&5.04&6&5\\
            8&5.98&9.89&9\\
            6.99&4.99&9&9.98
        \end{bmatrix}\begin{bmatrix}
            x_1''\\x_2''\\x_3''\\x_4''
        \end{bmatrix} = \begin{bmatrix}
            32\\23\\33\\31
        \end{bmatrix} \implies \begin{bmatrix}
            x_1''\\x_2''\\x_3''\\x_4''
        \end{bmatrix} = \begin{bmatrix}
            -81\\137\\-34\\22
        \end{bmatrix}.
    \]
    可见数据的微小变化会导致解的巨大变化，这是因为系数矩阵的条件数$\cond(A) = 32825/11$很大。
\end{example}

\begin{definition}
    {条件数}{condition number}
    给定诱导的矩阵范数$\norm \cdot$，可逆矩阵$A$的条件数(condition number)为
    \begin{equation}
        \cond(A) \equiv \norm A\norm{A\inv}.
    \end{equation}
\end{definition}

\begin{corollary}
    条件数的性质：
    \begin{itemize}
        \item $\cond(A) \geq 1$；
        \item $\cond(A\inv) = \cond(A)$；
        \item 对非零标量$c$，$\cond(cA) = \cond(A)$；
        \item 若$U$为正交矩阵，则$\cond_2(U) = 1$，且
        \begin{equation}
            \cond_2(A) = \cond_2(AU) = \cond_2(UA);
        \end{equation}
        \item 若$\lambda_1,\lambda_n$是$A$模最大与最小的特征值，则
        \begin{equation}
            \cond(A) \geq \frac{\abs{\lambda_1}}{\abs{\lambda_n}},
        \end{equation}
        若$A$对称，则$\cond_2(A) = \abs{\lambda_1}/\abs{\lambda_n}$；
        \item 由范数的等价性，可知条件数的等价性：
        \begin{subequations}
            \begin{alignat}{4}
                \frac1n&\cond_2&(A) &\leq \cond_1&(A) &\leq n\cond_2&(A),\\
                \frac1n&\cond_\infty&(A) &\leq \cond_2&(A) &\leq n\cond_\infty&(A),\\
                \frac1{n^2}&\cond_1&(A) &\leq \cond_\infty&(A) &\leq n^2\cond_1&(A).
            \end{alignat}
        \end{subequations}
    \end{itemize}
\end{corollary}

\begin{theorem}
    {解的扰动定理}{}
    给定可逆矩阵$A$和微小扰动$\Delta A$，满足
    \begin{equation}
        \label{eqn:perturbation condition}
        \frac{\norm{\Delta A}}{\norm A} < \frac1{\cond(A)},
    \end{equation}
    则$(A + \Delta A)$也可逆，考察线性方程组$Ax = b$及其扰动方程组
    \begin{equation}\label{eq:na-perturbed-linear-system}
        (A + \Delta A)(x + \Delta x) = b + \Delta b,
    \end{equation}
    当$b \neq 0$时有
    \begin{equation}
        \frac{\norm{\Delta x}}{\norm x} \leq \frac{\cond(A)}{1 - \norm{A\inv}\norm{\Delta A}}\biggp{\frac{\norm{\Delta A}}{\norm A} + \frac{\norm{\Delta b}}{\norm b}}.
    \end{equation}
\end{theorem}

\begin{proof}
    \eqref{eqn:perturbation condition}即$\norm{A\inv}\norm{\Delta A} < 1$。
    若$(A + \Delta A)$不可逆，即存在$x \neq 0$使得
    \[
        (A + \Delta A)x = 0 \iff A\inv\Delta A x = -x,
    \]
    即$A\inv\Delta A$有一个特征值$-1$，谱半径$\rho(A\inv\Delta A) \geq 1 > \norm{A\inv}\norm{\Delta A} \geq \norm{A\inv\Delta A}$与\eqref{eqn:rho<=norm}矛盾，故$(A + \Delta A)$可逆。令$B = (A + \Delta A)\inv$，则
    \begin{align*}
        (A + \Delta A)B = I \iff&B = A\inv - A\inv\Delta AB = A\inv(I - \Delta AB)\\
        \implies&\norm B \leq \norm{A\inv}\norm{I - \Delta AB} \leq \norm{A\inv}\Bigp{1 + \norm{\Delta A}\norm B},
    \end{align*}
    移项即得
    \begin{equation}
        \norm{(A + \Delta A)\inv} \leq \frac{\norm{A\inv}}{1 - \norm{A\inv}\norm{\Delta A}}.
    \end{equation}
    由
    \begin{align*}
        \Delta x &= (A + \Delta A)\inv(b + \Delta b) - x\\
        &= (A + \Delta A)\inv(b + \Delta b - (A + \Delta A)x)\\
        &= (A + \Delta A)\inv(\Delta b - \Delta Ax),
    \end{align*}
    两边取范数
    \begin{align*}
        \norm{\Delta x} &\leq \norm{(A + \Delta A)\inv}(\norm{\Delta b} + \norm{\Delta A}\norm x)\\
        &\leq \frac{\norm{A\inv}}{1 - \norm{A\inv}\norm{\Delta A}}\biggp{\frac{\norm{\Delta A}}{\norm A}\norm A\norm x + \frac{\norm{\Delta b}}{\norm b}\norm A\norm x}.
        \qedhere
    \end{align*}
\end{proof}

\begin{remark}
    因此条件数可以看成扰动方程组相对误差的放大倍数。
\end{remark}

\begin{theorem}
    {矩阵相对奇异性的度量}{relative singularity}
    若$A$可逆，定义所有使得$(A + \Delta A)$不可逆的$\Delta A$构成集合$S$，则
    \begin{equation}
        \min_{\Delta A\in S}\frac{\norm{\Delta A}_2}{\norm A_2} = \frac1{\cond_2(A)},
    \end{equation}
\end{theorem}

\begin{proof}
    当$\norm{A\inv}\norm{\Delta A} < 1$时，$(A + \Delta A)$可逆，故
    \[
        \min_{\Delta A\in S}\norm{\Delta A}_2 \geq \frac1{\norm{A\inv}_2}.
    \]
    由$\norm \cdot_2$的定义，$\exists x$且$\norm x_2 = 1$使得$\norm{A\inv x}_2 = \norm{A\inv}_2$，令$y = A\inv x/\norm{A\inv}_2$，并取
    \[
        \Delta A = -\frac{xy\transp}{\norm{A\inv}_2},
    \]
    则$\norm y_2 = 1$且
    \[
        (A + \Delta A)y = \frac{x}{\norm{A\inv}_2} - \frac{xy\transp y}{\norm{A\inv}_2} = 0,
    \]
    故$(A + \Delta A)$不可逆，又
    \[
        \norm{\Delta A}_2 = \max_{\norm z_2=1}\norm{\Delta Az}_2 = \frac{\norm x_2}{\norm{A\inv}_2}\max_{\norm z_2=1}\abs{y\transp z} = \frac1{\norm{A\inv}_2}.
        \qedhere
    \]
\end{proof}

\begin{remark}
    因此可逆矩阵到最接近的奇异矩阵的相对距离在2-范数意义下就是2-条件数的倒数。当条件数很大时，矩阵与奇异矩阵的相对距离很小，称为病态(ill conditioned)。
\end{remark}

\begin{theorem}
    {近似解的相对误差}{}
    若$x,x'$分别是方程组$Ax = b$的精确解和近似解，$x'$的剩余$r = b - Ax'$，则
    \begin{equation}
        \frac1{\cond(A)}\frac{\norm r}{\norm b} \leq \frac{\norm{x' - x}}{\norm x} \leq \cond(A)\frac{\norm r}{\norm b}.
    \end{equation}
\end{theorem}

\begin{proof}
    由$A(x' - x) = -r$和$\norm A\norm x \geq \norm b$可得
    \[
        \norm r \leq \norm A\norm{x' - x} = \frac{\cond(A)}{\norm{A\inv}}\norm{x' - x},
    \]
    两边除$\norm b$，由$x = A\inv b$可得
    \[
        \frac{\norm r}{\norm b} \leq \cond(A)\frac{\norm{x' - x}}{\norm{A\inv}\norm b} \leq \cond(A)\frac{\norm{x' - x}}{\norm x};
    \]
    另一方面，由$x' - x = -A\inv r$可得
    \[
        \norm{x' - x} = \norm{A\inv r} \leq \norm{A\inv}\norm r = \cond(A)\frac{\norm r}{\norm A}.
    \]
    两边同除$\norm x$，由$Ax = b$可得
    \[
        \frac{\norm{x' - x}}{\norm x} \leq \cond(A)\frac{\norm r}{\norm A\norm x} \leq \cond(A)\frac{\norm r}{\norm b}.
    \]
    综上，两边不等式均得证。
\end{proof}

\begin{remark}
    这说明当方程组病态时，即使剩余$\norm r$比较小，解的相对误差仍可能很大。
\end{remark}

\begin{example}
    {Hilbert矩阵}{Hilbert matrix}
    Hilbert矩阵$H_{ij} := 1/(i + j - 1)$
    \begin{equation*}
        H_n := \begin{bmatrix}
            1&1/2&\cdots&1/(n + 1)\\
            1/2&1/3&\cdots&1/(n + 2)\\
            \vdots&\vdots&\ddots&\vdots\\
            1/(n + 1)&1/(n + 2)&\cdots&1/(2n + 1)
        \end{bmatrix}
    \end{equation*}
    的条件数增长很快：
    \begin{equation}\label{eq:na-hilbert-condition-asymptotic}
        \cond_2(H_n) = \bigo\biggp{\frac{(1 + \sqrt2)^{4n}}{\sqrt n}}.
    \end{equation}
\end{example}

\begin{theorem}
    {条件数与数值精度}{}
    用直接法解方程组$Ax = b$，$A,b$的元素有效位数为$s$而$\cond(A)$的数量级为$t$，则求得$x$分量有效位数约为$s - t$。
\end{theorem}

\paragraph{病态方程组的解法}

除采用更高精度的运算外，另一个更有效的方法是对原方程进行预处理：
\begin{equation}\label{eq:na-two-sided-preconditioning}
    Ax = b \iff PAQ(Q\inv x) = Pb,
\end{equation}
从而降低系数矩阵的条件数：$\cond(PAQ) \ll \cond(A)$。一般$P,Q$可选择为三角矩阵或对角矩阵。

\begin{example}
    {预处理例子}{}
    方程组
    \[
        \begin{bmatrix}
            10&10^5\\1&1
        \end{bmatrix}\begin{bmatrix}
            x_1\\x_2
        \end{bmatrix} = \begin{bmatrix}
            10^5\\2
        \end{bmatrix} \implies \begin{bmatrix}
            x_1\\x_2
        \end{bmatrix} = \frac1{9999}\begin{bmatrix}
            10000\\9998
        \end{bmatrix}.
    \]
    系数矩阵的条件数$\cond_2(A) = 100010$很大，左乘$D = \diag(10^{-5},1)$平衡：
    \[
        \begin{bmatrix}
            10^{-4}&1\\1&1
        \end{bmatrix}\begin{bmatrix}
            x_1\\x_2
        \end{bmatrix} = \begin{bmatrix}
            1\\2
        \end{bmatrix},
    \]
    系数矩阵的条件数$\cond_2(DA) = 940/359$得到了有效降低。
\end{example}

\section{迭代法}

除了直接解法(如Gauss消元法)外，线性方程组还可以用迭代法(iterative method)求解。
所谓迭代法，是指通过算法生成一个向量序列$x^{(1)},x^{(2)},\ldots$，使其尽快收敛到方程$Ax = b$的解$x^*$。

\subsection{向量序列和矩阵序列的极限}

\begin{definition}
    {向量序列的极限}{limit of vector array}
    给定赋范空间$(\FF^n,\norm \cdot)$以及一组向量序列$x^{(1)},x^{(2)},\ldots$，若$\exists x \in \FF^n$使得
    \begin{equation}
        \lim_{i\to\infty}\norm{x^{(i)} - x} = 0 \iff \lim_{i\to\infty}x^{(i)} = x.
    \end{equation}
    则称$x^{(i)}$收敛于$x$。
\end{definition}

\begin{remark}
    由于有限维线性空间中所有范数都是等价的，向量序列的收敛性不依赖于范数的选择。
\end{remark}

\begin{theorem}
    {向量序列的收敛性判定}{}
    向量序列收敛等价于其各分量序列收敛：
    \begin{equation}
        \lim_{i\to\infty}x^{(i)} = x \iff \lim_{i\to\infty}x^{(i)}_j = x_j,\enspace \forall j = 1,\ldots,n.
    \end{equation}
\end{theorem}

\begin{proof}
    选取无穷范数：
    \[
        \lim_{i\to\infty}\norm{x^{(i)} - x}_\infty = 0 \iff \lim_{i\to\infty}\max_j\abs{x^{(i)}_j - x_j} = 0 \iff \lim_{i\to\infty}\abs{x^{(i)}_j - x_j} = 0,\enspace\forall j.
        \qedhere
    \]
\end{proof}

\begin{remark}
    类似地，我们可以定义矩阵序列的收敛性。
\end{remark}

\begin{theorem}
    {收敛到零矩阵}{limit=O}
    矩阵序列收敛到零矩阵的充要条件是它作用在每个向量上都趋于零：
    \begin{equation}\label{eq:na-matrix-sequence-convergence}
        \lim_{k\to\infty}A^{(k)} = O
        \iff \lim_{k\to\infty}A^{(k)}x = 0\quad\text{对每个 }x \in \FF^n.
    \end{equation}
\end{theorem}

\begin{proof}
    必要性：对于任一矩阵的从属范数，
    \[
        \norm{A^{(k)}x} \leq \norm{A^{(k)}}\norm x \to 0;
    \]
    充分性：其矩阵各分量
    \[
        \abs{a_{ij}^{(k)}} = \abs{e_i\transp A^{(k)}e_j} \leq \norm{A^{(k)}e_j}_\infty \to 0.
        \qedhere
    \]
\end{proof}

\begin{theorem}
    {矩阵幂序列的收敛性}{matrix exponent convergence}
    给定$B \in \FF^{n\times n}$，下面三个命题等价：
    \begin{enumerate}
        \item $\textstyle \lim_{k\to\infty}B^k = O$；
        \item $\rho(B) < 1$；
        \item 存在一个矩阵范数$\norm \cdot$使得$\norm B < 1$。
    \end{enumerate}
\end{theorem}

\begin{proof}
    $(1) \implies (2)$：
    采取反证法，若$\lambda$为$B$的特征值且$\abs\lambda \geq 1$其对应特征向量为$x$，则
    \[
        \norm{B^kx} = \abs\lambda^k\norm x\not \to 0;
    \]

    $(2) \implies (3)$：
    由$\rho(B) < 1$，给定$\epsilon = (1 - \rho(B))/2 > 0$，
    由\thmref{thm:norm<=rho+eps}，
    存在一个矩阵范数使得$\norm B \leq \rho(B) + \epsilon < 1$；

    $(3) \implies (1)$：
    由$\norm{B^k} \leq \norm B^k \to 0$即得。
\end{proof}

\begin{theorem}
    {矩阵幂的矩阵范数与谱半径}{k-root norm B^k=rho}
    给定$B \in \FF^{n\times n}$和任一矩阵范数$\norm \cdot$，均有
    \begin{equation}
        \lim_{k\to\infty}\norm{B^k}^{1/k} = \rho(B).
    \end{equation}
\end{theorem}

\begin{proof}
    采用夹逼定理。
    一方面，由\thmref{thm:rho<=norm}，
    \[
        \rho(B) = \rho(B^k)^{1/k} \leq \norm{B^k}^{1/k};
    \]
    另一方面，$\forall\epsilon > 0$，记
    \[
        B_\epsilon = \frac{B}{\rho(B) + \epsilon},
    \]
    则$\rho(B_\epsilon) < 1$，由\thmref{thm:matrix exponent convergence}，$\textstyle \lim_{k\to\infty}B_\epsilon^k = O$，即$\exists N = N(\epsilon) > 0$使得$\forall k > N$，
    \[
        \norm{B_\epsilon^k} = \frac{\norm{B^k}}{(\rho(B) + \epsilon)^k} < 1 \iff \norm{B^k}^{1/k} < \rho(B) + \epsilon \iff \lim_{k\to\infty}\norm{B^k}^{1/k} \leq \rho(B).
    \]
    综上，极限相等。
\end{proof}

\subsection{迭代公式的构造}

若线性方程$Ax = b$与方程$x = Bx + f$同解，则可以构造迭代公式：
\begin{equation}
    x^{(k+1)} = Bx^{(k)} + f,
\end{equation}
其中$B$称为迭代矩阵。若$x^{(k)}$收敛，则收敛值$x^*$就是方程的解。

\begin{theorem}
    {迭代法的收敛性分析}{convergence of iterative method}
    迭代法$x^{(k+1)} = Bx^{(k)} + f$收敛与以下命题等价：
    \begin{enumerate}
        \item $\textstyle \lim_{k\to\infty}B^k = O$；
        \item $\rho(B) < 1$；
        \item 存在一种矩阵范数$\norm \cdot$使得$\norm B < 1$。
    \end{enumerate}
\end{theorem}

\begin{proof}
    记误差向量$e^{(k)} := x^{(k)} - x^*$，则$e^{(k+1)} = Be^{(k)}$。迭代法收敛意味着$\forall e^{(0)} \in \FF^n$均有
    \[
        \lim_{k\to\infty}e^{(k)} = \lim_{k\to\infty}B^ke^{(0)} = 0,
    \]
    由\thmref{thm:limit=O}，$\textstyle \lim_{k\to\infty}B^k = O$，由\thmref{thm:matrix exponent convergence} 知三个命题等价。
\end{proof}

\begin{theorem}
    {迭代法的误差分析}{error of iterative method}
    若迭代法收敛，第$k$次迭代的误差满足
    \begin{equation}
        \norm{e^{(k)}} \leq \frac{\norm B}{1 - \norm B}\norm{x^{(k)} - x^{(k-1)}} \leq \frac{\norm B^k}{1 - \norm B}\norm{x^{(1)} - x^{(0)}}.
    \end{equation}
\end{theorem}

\begin{proof}
    由$e^{(k)} = Be^{(k-1)} = B(e^{(k-1)} - e^{(k)}) + Be^{(k)}$，故
    \[
        \norm{e^{(k)}} \leq \norm{B(e^{(k-1)} - e^{(k)})} + \norm{Be^{(k)}} \leq \norm B\Bigp{\norm{e^{(k-1)} - e^{(k)}} + \norm{e^{(k)}}},
    \]
    移项即得第一个不等号；又$e^{(k)} - e^{(k-1)} = B(e^{(k-1)} - e^{(k-2)})$可得
    \[
        \norm{e^{(k)} - e^{(k-1)}} \leq \norm B\norm{e^{(k-1)} - e^{(k-2)}},
    \]
    反复迭代即得第二个不等号。
\end{proof}

\begin{theorem}
    {迭代法的收敛速度}{converging speed of iterative method}
    若给定$\epsilon > 0$，希望
    \[
        \frac{\norm{e^{(k)}}}{\norm{e^{(0)}}} \leq \epsilon,
    \]
    则要求迭代次数
    \begin{equation}
        \label{eqn:k<-lneps/R}
        k \geq \frac{-\ln\epsilon}{-\ln\norm{B^k}^{1/k}}.
    \end{equation}
\end{theorem}

\begin{proof}
    按矩阵从属范数的定义
    \[
        \norm{B^k} = \max_{e^{(0)}\neq 0}\frac{\norm{B^ke^{(0)}}}{\norm{e^{(0)}}} = \max_{e^{(0)}\neq 0}\frac{\norm{e^{(k)}}}{\norm{e^{(0)}}}.
    \]
    因此只需$\norm{B^k} \leq \epsilon$即可。
\end{proof}

\begin{remark}
    式\eqref{eqn:k<-lneps/R}并不好处理，可以用近似：
    \begin{equation}\label{eq:na-asymptotic-iteration-count}
        k \approx \frac{-\ln\epsilon}{-\ln\rho(B)}.
    \end{equation}
\end{remark}

\begin{definition}
    {平均收敛率}{average converging speed}
    迭代法的平均收敛率为
    \begin{equation}
        R_k(B) := -\ln\norm{B^k}^{1/k}.
    \end{equation}
    其渐进收敛率$R(B) := R_\infty(B) = -\ln\rho(B)$。
\end{definition}

% \section{(单步)定常线性迭代}

\subsection{Richardson迭代法}

\begin{example}
    {Richardson迭代法}{Richardson iterative method}
    最简单的情形$x = (I - A)x + b$，称为Richardson迭代：
    \begin{equation}
        x^{(k+1)} = (I - A)x^{(k)} + b,
    \end{equation}
    收敛条件为$\rho(I - A) < 1$。
    定义残差$r^{(k)} := b - Ax^{(k)}$，则有递推关系：
    \begin{subequations}\label{eq:na-richardson-residual-update}
        \begin{align}
            r^{(k+1)} &= (I - A)r^{(k)},\\
            x^{(k+1)} &= x^{(k)} + r^{(k)},
        \end{align}
    \end{subequations}
    继而
    \[
        x^{(k)} = x^{(0)} + r^{(0)} + \cdots + r^{(k-1)} = x_0 + \sum_{i=0}^{k-1}(I - A)^ir^{(0)}.
    \]
\end{example}

\iffalse
\begin{definition}
    {Krylov子空间}{Krylov subspace}
    给定非零向量$v$和非零方阵$A$，定义
    \[
        \mathcal K_n(v):=\Span(v,Av,\ldots,A^{n-1}v),
    \]
    称为Krylov子空间。
\end{definition}

\begin{corollary}
    显然
    \[
        \mathcal K_1(v)\subset\mathcal K_2(v)\subset\cdots\subset\mathcal K_n(v)\subset\cdots
    \]
    最大维度$\sup_{n\to\infty}\dim\mathcal K_n(v)$称为$v$关于$A$的度。
\end{corollary}
\fi

\begin{definition}
    {预处理矩阵}{preconditioning matrix}
    若$M$可逆且$M\inv$是$A\inv$的一个较好的近似，对$M\inv Ax = M\inv b$使用Richardson迭代：
    \begin{equation}
        x^{(k+1)} = (I - M\inv A)x^{(k)} + M\inv b = x^{(k)} + M\inv r^{(k)},
    \end{equation}
    矩阵$M$称为预处理矩阵(preconditioning matrix)。
\end{definition}

\begin{remark}
    预处理矩阵的选择依据：
    方便计算、
    好近似。
\end{remark}

\subsection{Jacobi迭代法和Gauss--Seidel迭代法}

将$A$拆解为$A = D + L + U$的形式，其中$D$是对角部分，$L$是严格下三角部分，$U$是严格上三角部分。
\footnote{本节对$L,U$的定义与教材上的定义相差一个负号。}

\begin{example}
    {Jacobi迭代法}{Jacobi iterative method}
    预处理矩阵$M = D$，迭代矩阵$B_\text J = -D\inv(L + U)$，迭代的分量形式为：
    \begin{equation}
        x_k^{(i+1)} = \frac1{a_{kk}}\biggp{b_k - \sum_{j\neq k}a_{kj}x_j^{(i)}}.
    \end{equation}
\end{example}

\begin{example}
    {Gauss--Seidel迭代法}{GS iterative method}
    预处理矩阵$M = D + L$，迭代矩阵$B_\text G = -(D + L)\inv U$，迭代的分量形式为：
    \begin{equation}
        x_k^{(i+1)} = \frac1{a_{kk}}\biggp{b_k - \sum_{j=1}^{k-1}a_{kj}x_j^{(i+1)} - \sum_{j=k+1}^na_{kj}x_j^{(i)}}.
    \end{equation}
\end{example}

\begin{remark}
    Jacobi法和G--S法收敛$\iff \rho(B_\text J) < 1$和$\rho(B_\text G) < 1$。
    另外，有时容易验证矩阵$A$满足收敛的充分条件。
\end{remark}

\begin{theorem}
    {Jacobi法和G--S法收敛的充分条件：对角占优}{}
    若$A$各行均严格对角占优，或$A$不可约、各行弱对角占优且至少一行严格占优，则Jacobi法和G--S法收敛。
\end{theorem}

\begin{proof}
    对于Jacobi法，若$B_\text J$有特征值$\lambda$且$\abs\lambda \geq 1$，则
    \[
        \det(\lambda I - B_\text J) = 0 \implies \det(D + \lambda\inv(L + U)) = 0,
    \]
    若$A$严格对角占优(或不可约弱对角占优)，则$D + \lambda\inv(L + U)$也为严格对角占优(或不可约弱对角占优)，矛盾！故$\rho(B_\text J) < 1$，Jacobi法收敛。

    对于G--S法，设$B_\text G$有特征值$\lambda$且$\abs\lambda \geq 1$，则
    \[
        \det(\lambda I - B_\text G) = 0 \implies \det(D + L + \lambda\inv U) = 0,
    \]
    若$A$严格对角占优(或不可约弱对角占优)，则$D + L + \lambda\inv U$也为严格对角占优(或不可约弱对角占优)，矛盾！故$\rho(B_\text G) < 1$，G--S法收敛。
\end{proof}

\begin{theorem}
    {对称矩阵的Jacobi法和G--S法收敛的充要条件}{}
    设$A = D + L + U$实对称且对角元均为正。Jacobi法收敛当且仅当$A$与$2D - A$均正定；G--S法收敛当且仅当$A$正定。
\end{theorem}

\subsection{超松弛迭代法}

\begin{example}
    {逐次超松弛迭代}{successive overrelaxation iterative method}
    预处理矩阵$M = D/\omega + L$，迭代矩阵$L_\omega = (D + \omega L)\inv[(1 - \omega)D - \omega U]$，
    \begin{equation}
        x_k^{(i+1)} = (1 - \omega)x_k^{(i)} + \frac\omega{a_{kk}}\biggp{b_k - \sum_{j=1}^{k-1}a_{kj}x_j^{(i+1)} - \sum_{j=k+1}^na_{kj}x_j^{(i)}}.
    \end{equation}
    称为逐次超松弛迭代法(successive overrelaxation, SOR)，其中$\omega$称为松弛因子。
\end{example}

\begin{theorem}
    {$\rho(L_\omega)$与$\omega$的关系}{}
    若$A$可逆且对角元非0，则
    \begin{equation}
        \rho(L_\omega) \geq \abs{1 - \omega}.
    \end{equation}
\end{theorem}

\begin{proof}
    设$L_\omega$的特征值是$\lambda_1,\ldots,\lambda_n$，由于$L_\omega = (D + \omega L)\inv[(1 - \omega)D - \omega U]$可以写成下三角矩阵和上三角矩阵的乘积：
    \[
        \lambda_1\cdots\lambda_n = \det(L_\omega) = \det(D\inv)\det((1 - \omega)D) = (1 - \omega)^n,
    \]
    故
    \[
        \rho(L_\omega) = \max_i\abs{\lambda_i} \geq \abs{\lambda_1\cdots\lambda_n}^{1/n} = \abs{1 - \omega}.
        \qedhere
    \]
\end{proof}

\begin{corollary}
    SOR收敛的必要条件是$0 < \omega < 2$。
\end{corollary}

\begin{theorem}
    {对称正定矩阵的SOR收敛的充要条件}{}
    若$A$对称正定，则$0 < \omega < 2 \iff$SOR收敛。
\end{theorem}

\begin{proof}
    只需证明充分性。
    设$\lambda$和$x$是$L_\omega$的任意一个特征值和对应的特征向量，则
    \[
        [(1 - \omega)D - \omega U]x = \lambda(D + \omega L)x,
    \]
    上式两边与$x$内积得到
    \[
        % (1-\omega)\inp{Dx}x-\omega\inp{Ux}x=\lambda(\inp{Dx}x+\omega\inp{Lx}x)\implies
        \lambda\bigl(\inp{Dx}x + \omega\inp{Lx}x\bigr) = (1 - \omega)\inp{Dx}x - \omega\inp{Ux}x.
    \]
    由$A$正定，$D$也正定，记$\inp{Dx}x =: p > 0$，$\inp{Lx}x = a + \i b$，由$A$对称，$L\transp = U$，
    \[
        \inp{Ux}x = \inp{L\transp x}x = \inp x{Lx} = \overline{\inp{Lx}x} = a - \i b,
    \]
    由$A$正定，$\inp{Ax}x = p + 2a > 0$。因$0 < \omega < 2$，分母的实部$p + \omega a > p(1 - \omega/2) > 0$。因此
    \[
        \abs\lambda^2 - 1 = \frac{[(1 - \omega)p - \omega a]^2 + \omega^2b^2}{(p + \omega a)^2 + \omega^2b^2} - 1 = -\frac{p\omega(2 - \omega)(p + 2a)}{(p + \omega a)^2 + \omega^2b^2}
    \]
    右侧为负，故$\abs\lambda^2 < 1$，从而$\rho(L_\omega) < 1$，SOR收敛。
\end{proof}

\paragraph{最佳松弛因子}

希望选取最优的$\omega$使得迭代矩阵谱半径最小：
\begin{equation}
    \omega^* = \argmin_{\omega}\rho(L_\omega).
\end{equation}

\begin{theorem}
    {最佳松弛因子}{optimal overrelaxation}
    若$A$为对称正定的(分块)三对角矩阵，则$\rho(B_\text G) = \rho(B_\text J)^2 < 1$且SOR的最佳松弛因子为
    \begin{equation}
        \omega^* = \frac2{1 + \sqrt{1 - \rho(B_\text G)}}.
    \end{equation}
    且$\rho(L_{\omega^*}) = \omega^* - 1$。
\end{theorem}

\begin{remark}
    对于对称正定三对角矩阵，G--S法收敛速度是Jacobi法的两倍，SOR的收敛速度最快。
\end{remark}

\begin{example}
    {Poisson方程的五点差分格式}{five-point difference scheme}
    讨论二元函数$u = u(x,y)$满足Poisson方程
    \begin{equation}\label{eq:na-poisson-equation}
        \pv[2]ux + \pv[2]uy = f(x,y),
    \end{equation}
    边界条件$u|_{\p\Omega} = 0$。采用差分法，在均匀网格点$(x_i,y_i)$附近作Taylor展开：
    \[
        u_{i\pm 1,j} = u_{ij} \pm h\biggeval{\pv ux}_{ij} + \frac{h^2}2\biggeval{\pv[2]ux}_{ij} \pm \frac{h^3}6\biggeval{\pv[3]ux}_{ij} + \bigo(h^4).
    \]
    可得
    \begin{equation}\label{eq:na-centered-second-difference}
        \biggeval{\pv[2]ux}_{ij} = \frac{u_{i+1,j} + u_{i-1,j} - 2u_{ij}}{h^2} + \bigo(h^2)
    \end{equation}
    同理可得$\p^2u/\p y^2$。代入原Poisson方程得到五点差分格式：
    \begin{equation}
        u_{i+1,j} + u_{i-1,j} + u_{i,j+1} + u_{i,j-1} - 4u_{ij} = h^2f_{ij}.
    \end{equation}
    若矩形区域的两个方向各有$m$、$n$个内点，网格间距同为$h$，则可以写成线性方程组$Ax = b$的形式，其中
    \begin{alignat*}{2}
        x &= \begin{bmatrix}u_1\\\vdots\\u_n\end{bmatrix},
        &\qquad u_i &= \begin{bmatrix}u_{1i}\\\vdots\\u_{mi}\end{bmatrix},\\
        b &= -h^2\begin{bmatrix}f_1\\\vdots\\f_n\end{bmatrix},
        & f_i &= \begin{bmatrix}f_{1i}\\\vdots\\f_{mi}\end{bmatrix}.
    \end{alignat*}
    矩阵$A$及其对角块为
    \begin{subequations}\label{eq:na-poisson-block-matrix}
        \begin{align}
            A &= \begin{bmatrix}
                A_1&-I\\-I&A_2&\ddots\\ &\ddots&\ddots&-I\\ &&-I&A_n
            \end{bmatrix},\\
            A_i &= \begin{bmatrix}
                4&-1\\-1&4&\ddots\\ &\ddots&\ddots&-1\\ &&-1&4
            \end{bmatrix}.
        \end{align}
    \end{subequations}
    其中$A$是$mn$阶方阵，$A_i$是$m$阶方阵。
    故$A$是分块三对角正定矩阵，$B_\text J$的$mn$个特征值为
    \begin{equation}\label{eq:na-poisson-jacobi-spectrum}
        \lambda_{ij} = \frac12\biggp{\cos\frac{i\pi}{m+1} + \cos\frac{j\pi}{n+1}},\quad i = 1,\ldots,m,\enspace j = 1,\ldots,n.
    \end{equation}
    对单位正方形取$m = n$、$h = 1/(n+1)$，则$\rho(B_\text J) = \cos(\pi h)$，$\rho(B_\text G) = \cos^2(\pi h)$，最佳松弛因子为
    \begin{equation}\label{eq:na-poisson-optimal-sor}
        \omega^* = \frac2{1 + \sin(\pi h)},
    \end{equation}
    因此，Jacobi法和SOR的收敛速度分别为
    \begin{subequations}\label{eq:na-poisson-convergence-rates}
        \begin{alignat}{3}
            R(B_\text J) &= -\ln(\cos(\pi h)) = \frac{\pi^2}2h^2 + \frac{\pi^4}{12}h^4 + \bigo(h^6),\\
            R(L_{\omega^*}) &= -2\ln(\cos(\pi h)) + 2\ln(1 + \sin(\pi h)) = 2\pi h + \frac{\pi^3}3h^3 + \bigo(h^5).
        \end{alignat}
    \end{subequations}
    可见Jacobi法和SOR的收敛速度差了一个数量级。
\end{example}

\section{对称正定线性方程组的迭代法}

本节主要介绍共轭梯度法，一种用于求解对称正定线性方程组的迭代法。

\begin{definition}
    {对称正定算子}{symmetric positive definite operator}
    Hilbert空间$V$中的对称正定算子$A:V \to V$满足：$\forall x,y \in V$，
    \begin{itemize}
        \item $\inp{Ax}y = \inp x{Ay} \equiv \inp xy_A$；
        \item $\inp{Ax}x \geq 0$且$\inp{Ax}x = 0 \iff x = 0$。
    \end{itemize}
    可定义范数$\norm x_A := \sqrt{\inp{Ax}x}$。
\end{definition}

\begin{corollary}
    $\RR^n$中的内积和$n$阶对称正定矩阵$M$一一对应：
    \begin{equation}\label{eq:na-spd-inner-product}
        \inp xy_M = \inp{Mx}y.
    \end{equation}
\end{corollary}

\begin{theorem}
    {Ritz变分问题}{}
    考虑算子方程
    \[
        Ax = b,
    \]
    考虑$\varphi:V \to \RR$
    \begin{equation}
        \varphi(x) := \frac12\norm x_A^2 - \inp bx.
    \end{equation}
    若算子$A$对称正定，则
    \begin{equation}
        Ax^* = b \iff x^* = \argmin_{x\in V}\varphi(x).
    \end{equation}
\end{theorem}

\begin{proof}
    由$b = Ax^*$，$\varphi(x^*) = -\norm{x^*}_A^2/2$。
    $\forall x \in V$
    \[
        \varphi(x) - \varphi(x^*) = \frac12(\norm x_A^2 - 2\inp{Ax^*}x + \norm{x^*}_A^2) = \frac12\norm{x - x^*}^2_A.
    \]
    故$\varphi(x) \geq \varphi(x^*)$且$\varphi(x) = \varphi(x^*) \iff x = x^*$。
\end{proof}

\begin{corollary}
    迭代法的选择策略为：给定$x^{(k)}$，确定$x^{(k+1)}$使得
    \[
        \varphi(x^{(k+1)}) < \varphi(x^{(k)}).
    \]
\end{corollary}

\subsection{子空间方法}
\label{ssec:subspace method}

\begin{theorem}
    {子空间搜索}{subspace search}
    设$x + \Phi$是一个仿射(affine)平面，则
    \begin{equation}\label{eq:na-subspace-minimization}
        y^* = \argmin_{y\in\Phi}\varphi(x + y) \iff
        % \forall y\in\Phi,\enspace\inp{b-A(x+y^*)}y=0\iff
        b - A(x + y^*) \perp \Phi.
    \end{equation}
\end{theorem}

\begin{proof}
    由于
    \[
        \varphi(x + y) - \varphi(x^*) = \frac12\norm{x + y - x^*}_A^2
    \]
    故对于$y^*$应有
    \[
        \forall y \in \Phi,\enspace\inp{x^* - (x + y^*)}y_A = 0 \iff
        x^* - (x + y^*) \perp_A\Phi.
        \qedhere
    \]
\end{proof}

\begin{theorem}
    {一维极小搜索方法}{1-D minimal search}
    给定$x^{(k)}$和搜索方向$p^{(k)} \neq 0$，考虑一维极小值问题
    \begin{equation}
        \label{eqn:alphak}
        \alpha_k = \argmin_{\alpha\in\RR}\varphi(x^{(k)} + \alpha p^{(k)}) = \frac{\inp{r^{(k)}}{p^{(k)}}}{\norm{p^{(k)}}_A^2}.
    \end{equation}
    并使
    \begin{equation}
        x^{(k+1)} = x^{(k)} + \alpha_kp^{(k)},
    \end{equation}
    若$\alpha_k \neq 0$，则$\varphi(x^{(k+1)}) < \varphi(x^{(k)})$，且$r^{(k+1)} \perp p^{(k)}$。
\end{theorem}

\begin{proof}
    残差$r^{(k)} = b - Ax^{(k)}$，则
    \[
        \dd\alpha\varphi(x^{(k)} + \alpha p^{(k)}) = -\inp{r^{(k)}}{p^{(k)}} + \alpha\norm{p^{(k)}}_A^2 = 0.
        \qedhere
    \]
\end{proof}

\subsection{最速下降法}
\label{ssec:steepest descent method}

\begin{theorem}
    {最速下降法}{steepest descent method}
    搜索方向选择$\varphi$在$x^{(k)}$处下降最快的方向：$p^{(k)} = r^{(k)}$，称为最速下降法(steepest descent method)：
    \begin{subequations}
        \begin{align}
            r^{(k)} &= b - Ax^{(k)},\\
            \alpha_k &= \frac{\norm{r^{(k)}}^2}{\norm{r^{(k)}}_A^2},\\
            x^{(k+1)} &= x^{(k)} + \alpha_kr^{(k)}.
        \end{align}
    \end{subequations}
    则$\varphi(x^{(k)})$单调下降，以$\varphi(x^*)$为下界，且$x^{(k)}$在$k \to \infty$时收敛到$x^*$。
\end{theorem}

\begin{corollary}
    相邻两次搜索方向是正交的：$\inp{r^{(k+1)}}{r^{(k)}} = 0$，且
    \begin{equation}\label{eq:na-steepest-residual-update}
        r^{(k+1)} = (I - \alpha_k A)r^{(k)}.
    \end{equation}
\end{corollary}

\begin{theorem}
    {最速下降法的误差收敛性}{SDM error}
    记误差$e^{(k)} := x^{(k)} - x^*$，则误差的$A$-范数满足
    \begin{equation}
        \label{eqn:SDM e^k}
        \norm{e^{(k)}}_A \leq \biggp{\frac{\cond_2(A) - 1}{\cond_2(A) + 1}}^k\norm{e^{(0)}}_A.
    \end{equation}
\end{theorem}

\begin{lemma}
    $\forall x \neq 0$，
    \begin{equation}
        \frac{\norm x_A^2\norm x_{A\inv}^2}{\norm x^4} \leq \frac{(\cond_2(A) + 1)^2}{4\cond_2(A)}.
    \end{equation}
\end{lemma}

\begin{proof}
    设$\lambda_{\max},\lambda_{\min}$分别是对称正定算子$A$的最大最小特征值。
    则$\forall\lambda_i$为$A$的特征值，满足
    \[
        \frac1{\lambda_i} \leq \frac1{\lambda_{\min}} + \frac1{\lambda_{\max}} - \frac{\lambda_i}{\lambda_{\min}\lambda_{\max}}.
    \]
    则$\forall x$且$\norm x = 1$，
    \begin{align*}
        \frac{\norm x_A^2\norm x_{A\inv}^2}{\norm x^4} &= \frac{\inp{Ax}x\inp{A\inv x}x}{\inp xx^2} = \sum_i\lambda_ix_i^2\sum_i\frac1{\lambda_i}x_i^2\\
        &\leq \sum_i\lambda_ix_i^2\biggp{\frac1{\lambda_{\min}} + \frac1{\lambda_{\max}} - \frac1{\lambda_{\min}\lambda_{\max}}\sum_i\lambda_ix_i^2} \leq \frac{(\lambda_{\max} + \lambda_{\min})^2}{4\lambda_{\max}\lambda_{\min}}.
    \end{align*}
    由$A$对称，$\cond_2(A) = \abs{\lambda_{\max}}/\abs{\lambda_{\min}}$，由$A$正定，$\lambda_{\min} > 0$。
\end{proof}

\begin{proof}
    由$Ae^{(k)} = -r^{(k)}$，有$\inp{e^{(k)}}{r^{(k)}}_A = -\norm{r^{(k)}}^2$，$\norm{e^{(k)}}_A = \norm{r^{(k)}}_{A\inv}$，故
    \begin{align*}
        \norm{e^{(k+1)}}_A^2 &= \norm{e^{(k)} + \alpha_kr^{(k)}}_A^2
        = \norm{e^{(k)}}_A^2 + \alpha_k^2\norm{r^{(k)}}_A^2 + 2\alpha_k\inp{e^{(k)}}{r^{(k)}}_A\\
        &= \norm{e^{(k)}}_A^2 + \alpha_k^2\norm{r^{(k)}}_A^2 - 2\alpha_k\norm{r^{(k)}}^2 =
        \norm{e^{(k)}}_A^2 - \frac{\norm{r^{(k)}}^4}{\norm{r^{(k)}}_A^2}\\
        &= \norm{e^{(k)}}_A^2\biggp{1 - \frac{\norm{r^{(k)}}^4}{\norm{r^{(k)}}_A^2\norm{r^{(k)}}_{A\inv}^2}} \leq \norm{e^{(k)}}_A^2\biggp{1 - \frac{4\cond_2(A)}{(\cond_2(A) + 1)^2}}\\
        &= \norm{e^{(k)}}_A^2\biggp{\frac{\cond_2(A) - 1}{\cond_2(A) + 1}}^2.
    \end{align*}
    迭代即可得到。
\end{proof}

\begin{remark}
    虽然最速下降法总是收敛的，但当$A$病态时，$\cond_2(A) \gg 1$，收敛速度非常慢；当$\norm{r^{(k)}}$很小时，计算不稳定。
\end{remark}

\subsection{共轭梯度法}
\label{ssec:conjugate gradient method}

仍然考虑对称正定矩阵$A$和一维极小搜索，但不再沿着正交的$r^{(0)},r^{(1)},\ldots$，寻找另一组$p^{(0)},p^{(1)},\ldots$
对于第$k$次搜索：
\begin{equation}\label{eq:na-conjugate-direction-line-search}
    \alpha_k = \argmin_\alpha\varphi(x^{(k)} + \alpha p^{(k)}) \implies
    x^{(k+1)} = x^{(k)} + \alpha_kp^{(k)},
\end{equation}
由于
\[
    x^{(k+1)} = x^{(0)} + \alpha_0p^{(0)} + \cdots + \alpha_kp^{(k)}.
\]
我们希望$p^{(k)}$的选取能使得确定$\alpha_k$后，$x^{(k+1)}$是$x^{(0)} + \Span(p^{(0)},\ldots,p^{(k)})$中最优的，
即一维极小搜索的结果与$(k + 1)$维搜索结果相同：
\[
    (\alpha_0,\ldots,\alpha_k) = \argmin_{(\alpha_0,\ldots,\alpha_k)}\varphi(x^{(0)} + \alpha_0p^{(0)} + \cdots + \alpha_kp^{(k)}).
\]

\begin{definition}
    {$A$-正交向量组}{A-orthogonal}
    若向量组$p_1,\ldots,p_n,\ldots$满足$\forall i \neq j$均有
    \begin{equation}
        \inp{p_i}{p_j}_A = 0 \iff p_i \perp_A p_j.
    \end{equation}
    则称为$A$-正交向量组，也称$A$-共轭向量组。
\end{definition}

\begin{corollary}
    给定一个线性无关的向量组，可按照Gram--Schmidt正交化得到对应的$A$-正交向量组。
\end{corollary}

\begin{theorem}
    {多维子空间搜索问题}{}
    记$y \in \Span(p^{(0)},\ldots,p^{(k-1)})$，$x = x^{(0)} + y + \alpha p^{(k)}$，
    \[
        \varphi(x) = \varphi(x^{(0)} + y) + \alpha\inp{y}{p^{(k)}}_A - \alpha\inp{r^{(0)}}{p^{(k)}} + \frac{\alpha^2}2\norm{p^{(k)}}^2_A.
    \]
    上式出现了交叉项$\inp{y}{p^{(k)}}_A$，为了简化问题，自然可考虑使交叉项$\inp{y}{p^{(k)}}_A = 0$，这要求$p^{(k)}$与$\Span(p^{(0)},\ldots,p^{(k-1)})$是$A$-正交的。%(见\dfnref{dfn:A-orthogonal})

    由此便分解成了两个极小问题：
    \[
        \min_x\varphi(x) = \min_y\varphi(x^{(0)} + y) + \min_\alpha\biggb{-\alpha\inp{r^{(0)}}{p^{(k)}} + \frac{\alpha^2}2\norm{p^{(k)}}^2_A}
    \]
    前者满足$x^{(0)} + y^* = x^{(k)}$，后者得到
    \begin{equation}
        \alpha_k = \frac{\inp{r^{(0)}}{p^{(k)}}}{\norm{p^{(k)}}_A^2} = \frac{\inp{r^{(k)}}{p^{(k)}}}{\norm{p^{(k)}}_A^2}.
    \end{equation}
    由于$\inp{r^{(0)}}{p^{(k)}} = \inp{r^{(k)}}{p^{(k)}}$，故形式其实与\eqref{eqn:alphak}相同。
    并且有
    \begin{equation}
        \label{eqn:r^k+1=r^k-alpha_kAp^k}
        r^{(k+1)} = r^{(k)} - \alpha_kAp^{(k)}.
    \end{equation}
\end{theorem}

\begin{corollary}
    由子空间搜索\thmref{thm:subspace search}，$r^{(k)} \perp \Span(p^{(0)},\ldots,p^{(k-1)})$。
\end{corollary}

\begin{remark}
    与$p^{(0)},\ldots,p^{(k-1)}$ $A$-正交的$p^{(k)}$取法并不唯一。
\end{remark}

\paragraph{共轭梯度法}

当$r^{(k)} \neq 0$时，$p^{(0)},\ldots,p^{(k-1)},r^{(k)}$线性无关，可通过Gram--Schmidt正交化得到$p^{(k)}$：
\begin{equation}
    \label{eqn:Gram-Schmidt p^k}
    p^{(k)} = r^{(k)} - \sum_{i=0}^{k-1}\frac{\inp{r^{(k)}}{p^{(i)}}_A}{\norm{p^{(i)}}_A^2}p^{(i)},
    \tag{*}
\end{equation}

\begin{theorem}
    {}{}
    $r^{(k)} \perp_A\Span(p^{(0)},\ldots,p^{(k-2)})$。
\end{theorem}

\begin{proof}
    $r^{(k)} \perp \Span(p^{(0)},\ldots,p^{(k-1)}) = \Span(r^{(0)},\ldots,r^{(k-1)})$，由式\eqref{eqn:r^k+1=r^k-alpha_kAp^k}
    \[
        Ap^{(i)} = \frac{r^{(i)} - r^{(i+1)}}{\alpha_i} \perp r^{(k)},\enspace i = 0,\ldots,k - 2.
        \qedhere
    \]
\end{proof}

\begin{corollary}
    $p^{(k)}$为$r^{(k)}$和$p^{(k-1)}$的线性组合，式\eqref{eqn:Gram-Schmidt p^k}变为
    \[
        p^{(k)} = r^{(k)} - \frac{\inp{r^{(k)}}{p^{(k-1)}}_A}{\norm{p^{(k-1)}}_A^2}p^{(k-1)} =: r^{(k)} + \beta_{k-1}p^{(k-1)}.
    \]
    可以简化系数$\beta_k$的计算：
    \begin{equation}\label{eq:na-conjugate-gradient-beta}
        \beta_k = -\frac{\inp{r^{(k+1)}}{p^{(k)}}_A}{\norm{p^{(k)}}_A^2} = \frac{\norm{r^{(k+1)}}^2}{\alpha_k\norm{p^{(k)}}_A^2} = \frac{\norm{r^{(k+1)}}^2}{\norm{r^{(k)}}^2}.
    \end{equation}
\end{corollary}

\begin{theorem}
    {共轭梯度法}{conjugate gradient method}
    共轭梯度法(conjugate gradient method)初值
    \begin{equation}
        p^{(0)} = r^{(0)} = b - Ax^{(0)},
    \end{equation}
    迭代算法
    \begin{subequations}
        \begin{align}
            % r^{(k)}&=b-Ax^{(k)},\\
            \alpha_k &= \frac{\norm{r^{(k)}}^2}{\norm{p^{(k)}}_A^2},\\
            x^{(k+1)} &= x^{(k)} + \alpha_kp^{(k)},\\
            r^{(k+1)} &= r^{(k)} - \alpha_kAp^{(k)};\\
            \beta_k &= \frac{\norm{r^{(k+1)}}^2}{\norm{r^{(k)}}^2},\\
            p^{(k+1)} &= r^{(k+1)} + \beta_kp^{(k)}.
        \end{align}
    \end{subequations}
    若$r^{(k)} = 0$，则$x^{(k)} = x^*$，计算终止。
\end{theorem}

\begin{corollary}
    由于$r^{(0)},\ldots,r^{(k)}$两两正交，理论上最多$n = \dim(V)$步便可以得到精确解。
\end{corollary}

\begin{theorem}
    {梯度下降法的误差收敛性}{}
    误差$e^{(k)} := x^{(k)} - x^*$的$A$-范数满足
    \begin{equation}
        \label{eqn:CG e^k}
        \norm{e^{(k)}}_A \leq 2\biggp{\frac{\sqrt{\cond_2(A)} - 1}{\sqrt{\cond_2(A)} + 1}}^k\norm{e^{(0)}}_A.
    \end{equation}
    其中$\lambda_{\max},\lambda_{\min}$分别是对称正定算子$A$的最大最小特征值。
\end{theorem}

\begin{proof}
    记$\alpha = \lambda_{\min}/\lambda_{\max}$，取
    \[
        p_k^*(x) = \Biggdivision{T_k\biggp{\frac{2x - (1 + \alpha)}{1 - \alpha}}}{T_k\biggp{-\frac{1 + \alpha}{1 - \alpha}}},
    \]
    其中$T_k$是Chebyshev多项式，由于$\abs t \geq 1$时，
    \[
        T_k(t) = \frac12\biggb{\Bigp{t + \sqrt{t^2 - 1}}^k + \Bigp{t - \sqrt{t^2 - 1}}^k}\gtrapprox \frac12\Bigp{\abs t + \sqrt{t^2 - 1}}^k.
    \]
    则$\forall x \in (\alpha,1)$，
    \[
        \norm{p_k^*}_\infty = \Biggdivision{1}{T_k\biggp{-\frac{1 + \alpha}{1 - \alpha}}} \leq 2\biggp{\frac{1 + \alpha}{1 - \alpha} - \sqrt{\Bigp{\frac{1 + \alpha}{1 - \alpha}}^2 - 1}}^k = 2\biggp{\frac{1 - \sqrt\alpha}{1 + \sqrt\alpha}}^k.
        \qedhere
    \]
\end{proof}

\paragraph{共轭梯度法的预处理}
% \label{ssec:preconditioning}

从误差估计\eqref{eqn:CG e^k}可以看出，共轭梯度法的收敛速度强烈依赖于$A$的条件数，当$A$病态时收敛得很慢。为了改善收敛速度，可以先进行预处理(preconditioning)，降低矩阵的条件数。

\begin{theorem}
    {预处理的共轭梯度算法}{preconditioning conjugate gradient}
    给定对称正定的预处理矩阵$M$，则
    \begin{equation}\label{eq:na-preconditioned-system}
        Ax = b \iff M\inv Ax = M\inv b,
    \end{equation}
    对预处理后的方程应用共轭梯度算法：$p^{(0)} = M\inv r^{(0)}$，
    \begin{subequations}
        \begin{align}
            % r^{(k)}&=b-Ax^{(k)},\\
            \alpha_k &= \frac{\norm{r^{(k)}}_{M\inv}^2}{\norm{p^{(k)}}_A^2},\\
            x^{(k+1)} &= x^{(k)} + \alpha_kp^{(k)},\\
            r^{(k+1)} &= r^{(k)} - \alpha_kAp^{(k)};\\
            \beta_k &= \frac{\norm{r^{(k+1)}}_{M\inv}^2}{\norm{r^{(k)}}_{M\inv}^2},\\
            p^{(k+1)} &= M\inv r^{(k+1)} + \beta_kp^{(k)}.
        \end{align}
    \end{subequations}
    当$r^{(k)} = 0$时计算终止。
\end{theorem}

\begin{corollary}
    其误差仍满足式\eqref{eqn:CG e^k}，但条件数改善为$\cond_2(M\inv A)$。
\end{corollary}

\begin{remark}
    预处理矩阵的选择依据：
    方便计算、
    好近似。
    \begin{itemize}
        \item 代数方法：分解$A = M + N$，其中$M = LL\transp$对称正定、$L$尽可能稀疏，而$N$尽可能小，称为$A$的不完全Cholesky分解。

        比如取为对角部分$M = D$，在$A$的对角元素相差较大时收敛速度会大大提高。
        \item 分析方法：谱等价，$\forall x$
        \begin{equation}\label{eq:na-spectral-equivalence}
            c_1\norm x_A^2 \leq \norm x_M^2 \leq c_2\norm x_A^2,
        \end{equation}
        则$\rho(M\inv A) \in [c_2\inv,c_1\inv]$。
    \end{itemize}
\end{remark}
